\documentclass{amsart}
% For dealing with references we use the comment environment
\usepackage{verbatim}
\newenvironment{reference}{\comment}{\endcomment}
%\newenvironment{reference}{}{}
\newenvironment{slogan}{\comment}{\endcomment}
\newenvironment{history}{\comment}{\endcomment}

% For commutative diagrams we use Xy-pic
\usepackage[all]{xy}

% We use 2cell for 2-commutative diagrams.
\xyoption{2cell}
\UseAllTwocells

% We use multicol for the list of chapters between chapters
\usepackage{multicol}

% This is generally recommended for better output
\usepackage{lmodern}
\usepackage[T1]{fontenc}

% For cross-file-references

% Package for hypertext links:
\usepackage{hyperref}

% For any local file, say "hello.tex" you want to link to please

% Theorem environments.
%
\theoremstyle{plain}
\newtheorem{theorem}[subsection]{Theorem}
\newtheorem{proposition}[subsection]{Proposition}
\newtheorem{lemma}[subsection]{Lemma}

\theoremstyle{definition}
\newtheorem{definition}[subsection]{Definition}
\newtheorem{example}[subsection]{Example}
\newtheorem{exercise}[subsection]{Exercise}
\newtheorem{situation}[subsection]{Situation}

\theoremstyle{remark}
\newtheorem{remark}[subsection]{Remark}
\newtheorem{remarks}[subsection]{Remarks}

\numberwithin{equation}{subsection}

% Macros
%
\def\lim{\mathop{\mathrm{lim}}\nolimits}
\def\colim{\mathop{\mathrm{colim}}\nolimits}
\def\Spec{\mathop{\mathrm{Spec}}}
\def\Hom{\mathop{\mathrm{Hom}}\nolimits}
\def\Ext{\mathop{\mathrm{Ext}}\nolimits}
\def\SheafHom{\mathop{\mathcal{H}\!\mathit{om}}\nolimits}
\def\SheafExt{\mathop{\mathcal{E}\!\mathit{xt}}\nolimits}
\def\Sch{\mathit{Sch}}
\def\Mor{\mathop{\mathrm{Mor}}\nolimits}
\def\Ob{\mathop{\mathrm{Ob}}\nolimits}
\def\Sh{\mathop{\mathit{Sh}}\nolimits}
\def\NL{\mathop{N\!L}\nolimits}
\def\CH{\mathop{\mathrm{CH}}\nolimits}
\def\proetale{{pro\text{-}\acute{e}tale}}
\def\etale{{\acute{e}tale}}
\def\QCoh{\mathit{QCoh}}
\def\Ker{\mathop{\mathrm{Ker}}}
\def\Im{\mathop{\mathrm{Im}}}
\def\Coker{\mathop{\mathrm{Coker}}}
\def\Coim{\mathop{\mathrm{Coim}}}

% Boxtimes
%
\DeclareMathSymbol{\boxtimes}{\mathbin}{AMSa}{"02}

%
% Macros for moduli stacks/spaces
%
\def\QCohstack{\mathcal{QC}\!\mathit{oh}}
\def\Cohstack{\mathcal{C}\!\mathit{oh}}
\def\Spacesstack{\mathcal{S}\!\mathit{paces}}
\def\Quotfunctor{\mathrm{Quot}}
\def\Hilbfunctor{\mathrm{Hilb}}
\def\Curvesstack{\mathcal{C}\!\mathit{urves}}
\def\Polarizedstack{\mathcal{P}\!\mathit{olarized}}
\def\Complexesstack{\mathcal{C}\!\mathit{omplexes}}
% \Pic is the operator that assigns to X its picard group, usage \Pic(X)
% \Picardstack_{X/B} denotes the Picard stack of X over B
% \Picardfunctor_{X/B} denotes the Picard functor of X over B
\def\Pic{\mathop{\mathrm{Pic}}\nolimits}
\def\Picardstack{\mathcal{P}\!\mathit{ic}}
\def\Picardfunctor{\mathrm{Pic}}
\def\Deformationcategory{\mathcal{D}\!\mathit{ef}}

\setlength{\emergencystretch}{3em}
\begin{document}
\title[Stacks Project, Tag 0E9Z]{263 --- Fundamental class of a local complete intersection and detection of smoothness}
\maketitle
\noindent Copyright (C) 2005--2025 Johan de Jong. Stacks Project authors. Source: \href{https://stacks.math.columbia.edu/tag/0E9Z}{Tag 0E9Z}.

\noindent Editorial difficulty note (not author text). Difficulty H2: The proof constructs a determinant map from the conormal differential sequence to the relative dualizing line. It proves that invertibility of this map forces the maximal-rank Jacobian condition for a local complete-intersection presentation. Both the fundamental class construction and the converse smoothness criterion are substantive obligations..

\noindent Editorial provenance note (not author text). History: excerpt from revision \texttt{a0444\allowbreak 6e57e\allowbreak c1fbc\allowbreak 252a8\allowbreak 71afc\allowbreak ec775\allowbreak 2fb28\allowbreak 07b14}, duality.tex, lines 7556--7625. Reference presentation was adapted; original-author context and core support blocks were retained or added. The mathematical wording is unchanged. Licensed under GNU FDL 1.2 or later; no invariant sections or cover texts. See ../licenses/COPYING. Context blocks reproduce author definitions and supporting results; they are not additional counted problems.

\begin{lemma}
\label{lemma-fundamental-class-lci}
Let $Y$ be a Noetherian scheme. Let $f : X \to Y$ be a
local complete intersection morphism which factors
as an immersion $X \to P$ followed by a proper smooth morphism $P \to Y$.
Let $r$ be the locally constant function on
$X$ such that $\omega_{X/Y} = H^{-r}(f^!\mathcal{O}_Y)$
is the unique nonzero cohomology sheaf of $f^!\mathcal{O}_Y$, see
Lemma \href{https://stacks.math.columbia.edu/tag/0B6V}{lemma lci shriek (Tag 0B6V)}.
Then there is a map
$$
\wedge^r\Omega_{X/Y} \longrightarrow \omega_{X/Y}
$$
which is an isomorphism on the stalk at a point $x$ if and only
if $f$ is smooth at $x$.
\end{lemma}

\begin{proof}
The assumption implies that $X$ is compactifiable over $Y$ hence $f^!$
is defined, see Section \href{https://stacks.math.columbia.edu/tag/0A9Y}{section upper shriek (Tag 0A9Y)}.
Let $j : W \to P$ be an open subscheme such that
$X \to P$ factors through a closed immersion $i : X \to W$.
Moreover, we have $f^! = i^! \circ j^! \circ g^!$ where
$g : P \to Y$ is the given morphism.
We have $g^!\mathcal{O}_Y = \wedge^d\Omega_{P/Y}[d]$ by
Lemma \href{https://stacks.math.columbia.edu/tag/0BRT}{lemma smooth proper (Tag 0BRT)} where $d$ is the locally
constant function giving the relative dimension of $P$ over $Y$.
We have $j^! = j^*$. We have $i^!\mathcal{O}_W = \wedge^c\mathcal{N}[-c]$
where $c$ is the codimension of $X$ in $W$ (a locally constant
function on $X$) and where $\mathcal{N}$ is the normal sheaf of
the Koszul-regular immersion $i$, see Lemma \href{https://stacks.math.columbia.edu/tag/0BR0}{lemma regular immersion (Tag 0BR0)}.
Combining the above we find
$$
f^!\mathcal{O}_Y =
\left(\wedge^c\mathcal{N} \otimes_{\mathcal{O}_X}
\wedge^d\Omega_{P/Y}|_X\right)[d - c]
$$
where we have also used Lemma \href{https://stacks.math.columbia.edu/tag/0B6U}{lemma perfect comparison shriek (Tag 0B6U)}.
Thus $r = d|_X - c$ as locally constant functions on $X$.
The conormal sheaf of $X \to P$ is the module
$\mathcal{I}/\mathcal{I}^2$ where $\mathcal{I} \subset \mathcal{O}_W$
is the ideal sheaf of $i$, see
Morphisms, Section \href{https://stacks.math.columbia.edu/tag/01R1}{section conormal sheaf (Tag 01R1)}.
Consider the canonical exact sequence
$$
\mathcal{I}/\mathcal{I}^2 \to
\Omega_{P/Y}|_X \to \Omega_{X/Y} \to 0
$$
of Morphisms, Lemma \href{https://stacks.math.columbia.edu/tag/01UZ}{lemma differentials relative immersion (Tag 01UZ)}.
We obtain our map by an application of Lemma \href{https://stacks.math.columbia.edu/tag/0E9Y}{lemma determinant (Tag 0E9Y)}.

\medskip\noindent
If $f$ is smooth at $x$, then the map is an isomorphism by an application of
Lemma \href{https://stacks.math.columbia.edu/tag/0E9Y}{lemma determinant (Tag 0E9Y)}
and the fact that $\Omega_{X/Y}$ is locally free at $x$
of rank $r$. Conversely, assume that our map is an isomorphism on stalks
at $x$. Then the lemma shows that $\Omega_{X/Y}$ is free of rank $r$
after replacing $X$ by an open neighbourhood of $x$.
On the other hand, we may also assume that $X = \Spec(A)$ and
$Y = \Spec(R)$ where $A = R[x_1, \ldots, x_n]/(f_1, \ldots, f_m)$
and where $f_1, \ldots, f_m$ is a Koszul regular sequence
(this follows from the definition of local complete intersection morphisms).
Clearly this implies $r = n - m$. We conclude that the rank of the matrix
of partials $\partial f_j/\partial x_i$ in the residue field at $x$ is $m$.
Thus after reordering the variables we may assume
the determinant of $(\partial f_j/\partial x_i)_{1 \leq i, j \leq m}$
is invertible in an open neighbourhood of $x$. It follows
that $R \to A$ is smooth at this point, see for example
Algebra, Example \href{https://stacks.math.columbia.edu/tag/00T8}{example make standard smooth (Tag 00T8)}.
\end{proof}
\end{document}
